\subsection{The GaGa Model for Differential Expression}
\label{sec:gaga-pig}

\cite{Rossell2009GaGa} introduced the Gamma-Gamma, or GaGa, model for
detecting differential expression in microarray experiments.  The model is
particularly relevant for small-sample genomic studies, where thousands of
genes are measured but each experimental condition has only a small number of
replicates.  In this setting, gene-specific likelihood fits are unstable, and
the GaGa model borrows information across genes through an empirical-Bayes
hierarchy.  A central computational difficulty is that posterior inference for
the gene-specific gamma shape parameter involves ratios of gamma functions.
The original implementation evaluates this layer using a Stirling-based
approximation.  Our P-IG augmentation provides a direct posterior sampler for
the same gamma-shape layer, allowing us to replace the approximation while
leaving the empirical-Bayes GaGa structure unchanged.

Let $x_{ij}$ denote the expression measurement for gene $i$ in sample $j$, and
let $g(j)$ be the experimental group of sample $j$.  A pattern
$z_i \in \{0,\ldots,H-1\}$ partitions the experimental groups into clusters
$c \in \mathcal{C}_{z_i}$; groups in the same cluster share a mean expression
level, whereas groups in different clusters have distinct means.  Conditional
on $z_i=h$, the GaGa model used here is
\begin{align}
x_{ij}\mid \alpha_i,\lambda_{ic}, z_i=h
  &\sim \Ga\left(\alpha_i,\alpha_i/\lambda_{ic}\right),
  \qquad c=c_h\{g(j)\},
  \label{eq:gaga-likelihood}
\\
q_{ic}=\lambda_{ic}^{-1}\mid a_0,\nu
  &\sim \Ga\left(a_0,a_0/\nu\right),
  \label{eq:gaga-lambda-prior}
\\
\alpha_i\mid b_\alpha,\nu_\alpha
  &\sim \Ga\left(b_\alpha,b_\alpha/\nu_\alpha\right),
  \qquad
z_i\mid \pi \sim \mathrm{Categorical}(\pi_0,\ldots,\pi_{H-1}).
  \label{eq:gaga-alpha-prior}
\end{align}
Here we use the rate parameterization of the gamma distribution.  The
hyperparameters
$\theta=(a_0,\nu,b_\alpha,\nu_\alpha,\pi)$ are estimated by the original GaGa
empirical-Bayes EM algorithm.  In the comparisons below we deliberately keep
these estimates fixed and replace only the posterior scoring step.  This
isolates the effect of replacing the Stirling approximation with P-IG
augmentation.

For a fixed pattern $h$ and cluster $c$, define
\[
n_{ic}=|\{j:g(j)\in c\}|,\qquad
S_{ic}=\sum_{j:g(j)\in c}x_{ij},\qquad
L_{ic}=\sum_{j:g(j)\in c}\log x_{ij},\qquad
r_0=a_0/\nu .
\]
After integrating out $q_{ic}$, the collapsed cluster likelihood is
\begin{align}
m_{ic}(\alpha_i)
&=
\frac{\alpha_i^{n_{ic}\alpha_i}}{\Gamma(\alpha_i)^{n_{ic}}}
\exp\{(\alpha_i-1)L_{ic}\}
\frac{r_0^{a_0}}{\Gamma(a_0)}
\frac{\Gamma(a_0+n_{ic}\alpha_i)}
     {(r_0+\alpha_i S_{ic})^{a_0+n_{ic}\alpha_i}} .
\label{eq:gaga-collapsed}
\end{align}
Thus, for fixed $\alpha_i$, the posterior probability of pattern $h$ is
\begin{align}
\Pr(z_i=h\mid \alpha_i,x_i,\theta)
\propto
\pi_h \prod_{c\in\mathcal{C}_h}m_{ic}(\alpha_i).
\label{eq:gaga-pattern-update}
\end{align}
Given $z_i=h$, the collapsed conditional density of $\alpha_i$ is proportional
to
\begin{align}
p(\alpha_i\mid z_i=h,x_i,\theta)
\propto
\alpha_i^{b_\alpha-1}\exp\{-(b_\alpha/\nu_\alpha)\alpha_i\}
\prod_{c\in\mathcal{C}_h}m_{ic}(\alpha_i).
\label{eq:gaga-alpha-collapsed}
\end{align}
The difficult component is the gamma-shape ratio
$\prod_c \Gamma(a_0+n_{ic}\alpha_i)/\Gamma(\alpha_i)^{n_{ic}}$, which is the
part approximated in the original GaGa scoring algorithm.

Our sampler targets \eqref{eq:gaga-alpha-collapsed} using P-IG augmentation.
Let $N_i=\sum_c n_{ic}$ and draw
\[
\omega_{is}\mid \alpha_i \iid \PIG(\alpha_i),
\qquad s=1,\ldots,N_i,
\]
for the reciprocal gamma terms, and draw Damsleth auxiliary variables
\[
\tau_{ic}\mid \alpha_i \sim \Ga(a_0+n_{ic}\alpha_i,1),
\qquad c\in\mathcal{C}_{z_i},
\]
for the numerator gamma terms.  With
$\Omega_i=\sum_{s=1}^{N_i}\omega_{is}$ and
\[
B_i
=N_i\gamma+\sum_{c}L_{ic}-b_\alpha/\nu_\alpha
  +\sum_c n_{ic}\log\tau_{ic},
\]
the augmented conditional density for $y_i=\log\alpha_i$, including the
Jacobian, is
\begin{align}
\ell_i(y_i)
&=
(N_i+b_\alpha)y_i
-\Omega_i\exp(2y_i)
+B_i\exp(y_i)
+R_i\{\exp(y_i)\},
\label{eq:gaga-pig-logalpha}
\\
R_i(\alpha)
&=
N_i\alpha\log\alpha
-\sum_c(a_0+n_{ic}\alpha)\log(r_0+\alpha S_{ic}).
\nonumber
\end{align}
We sample $y_i$ by stepping-out slice sampling from
\eqref{eq:gaga-pig-logalpha}.  This is an exact Markov transition for the
P-IG-augmented conditional distribution.  The resulting per-gene sampler is:
\begin{enumerate}
\item sample $z_i$ from \eqref{eq:gaga-pattern-update};
\item conditional on $z_i$, draw $\omega_{is}$ and $\tau_{ic}$;
\item sample $y_i=\log\alpha_i$ from \eqref{eq:gaga-pig-logalpha};
\item after burn-in, estimate posterior pattern probabilities by the
Rao--Blackwellized average
\[
\widehat p_{ih}
=
\frac{1}{T}\sum_{t=1}^{T}
\Pr\{z_i=h\mid \alpha_i^{(t)},x_i,\widehat\theta\}.
\]
\end{enumerate}
The differential-expression score is $1-\widehat p_{i0}$, where pattern $0$
is the all-equal null pattern.

\textit{Oracle simulation.}
We first studied the posterior scoring layer in isolation.  Data were generated
from the GaGa model under fixed oracle hyperparameters, and each data set was
scored three ways: (i) high-accuracy one-dimensional quadrature over
$\alpha_i$, treated as the gold standard; (ii) the original GaGa
Stirling-based scorer; and (iii) the P-IG MCMC scorer above.  We used
six regimes spanning the paper-like high-shape setting, benign moderate
settings, low-shape/high-CV settings, and weak-signal settings.  Each regime
used 50 replicated data sets with 400 genes.  The P-IG sampler used 1000
iterations, 400 burn-in iterations, thinning by 2, and P-IG truncation 50.

Table~\ref{tab:gaga-oracle-simulation} summarizes the results.  In the
paper-like and benign regimes, the original Stirling approximation is already
essentially indistinguishable from exact quadrature.  In these regimes P-IG
MCMC produces the same decisions up to Monte Carlo error.  The advantage of
the P-IG construction appears in the low-shape regimes.  In the severe stress
setting $(a_0=1.5,b_\alpha=0.8,\nu_\alpha=0.7,n_g=3)$, the mean absolute error
of the original Stirling posterior null probability is 0.0428, compared with
0.0014 for P-IG MCMC.  At a nominal posterior FDR threshold of 0.05, the exact
scorer selects on average 0.10 genes and P-IG MCMC selects 0.16 genes, whereas
the Stirling scorer selects 16.50 genes with empirical FDR 0.832.

\begin{table}[t]
\centering
\caption{Oracle GaGa simulation comparing original Stirling scoring and
P-IG MCMC scoring against exact quadrature.  MAE is the mean absolute error in
the posterior null probability $p_{i0}$ relative to exact quadrature.  FDR is
the empirical false discovery rate among genes selected at posterior FDR 0.05.
Numbers in parentheses are Monte Carlo standard errors over 50 replicates.}
\label{tab:gaga-oracle-simulation}
\smallskip
\renewcommand{\arraystretch}{1.12}
\setlength{\tabcolsep}{4pt}
{\footnotesize
\begin{tabular}{@{}llrrrl@{}}
\toprule
\textbf{Regime} & \textbf{Method} &
\textbf{MAE($p_{i0}$)} &
\textbf{Spearman} &
\textbf{Selected} &
\textbf{FDR} \\
\midrule
Paper-like & Stirling & 0.0000 (0.0000) & 1.0000 & 35.54 (0.87) & 0.038 (0.003) \\
           & P-IG MCMC & 0.0058 (0.0002) & 0.9983 & 33.42 (0.84) & 0.016 (0.003) \\
Benign & Stirling & 0.0000 (0.0000) & 1.0000 & 25.06 (0.74) & 0.051 (0.005) \\
       & P-IG MCMC & 0.0064 (0.0001) & 0.9944 & 25.14 (0.77) & 0.059 (0.005) \\
Moderate small-$n$ & Stirling & 0.0002 (0.0000) & 0.9998 & 2.86 (0.29) & 0.048 (0.025) \\
                  & P-IG MCMC & 0.0054 (0.0001) & 0.9975 & 3.08 (0.30) & 0.051 (0.023) \\
Low shape & Stirling & 0.0029 (0.0002) & 0.9886 & 0.40 (0.09) & 0.190 (0.055) \\
          & P-IG MCMC & 0.0017 (0.0000) & 0.9995 & 0.08 (0.04) & 0.000 (0.000) \\
Severe stress & Stirling & 0.0428 (0.0013) & 0.9048 & 16.50 (0.67) & 0.832 (0.012) \\
              & P-IG MCMC & 0.0014 (0.0000) & 0.9995 & 0.16 (0.05) & 0.020 (0.020) \\
Weak signal & Stirling & 0.0007 (0.0000) & 0.9918 & 0.00 (0.00) & 0.000 (0.000) \\
            & P-IG MCMC & 0.0010 (0.0000) & 0.9985 & 0.00 (0.00) & 0.000 (0.000) \\
\bottomrule
\end{tabular}}
\end{table}

\textit{Empirical comparison with GaGa.}
We next applied the same plug-in empirical-Bayes strategy to the two empirical
examples considered by \cite{Rossell2009GaGa}.\footnote{The empirical data used
here are the closest public versions we could recover, not the exact
preprocessed analysis objects used by \cite{Rossell2009GaGa}.  For the
Armstrong leukemia example, the original raw CEL archives needed to reproduce
the authors' \texttt{just.gcrma} preprocessing were not available from the
original links at the time of our replication.  We therefore used the processed
Schliep mirror, downloaded as \texttt{armstrong-2002-v2\_database.txt} from the
Schliep supplementary data page \texttt{armstrong-2002-v2.html}; we also checked
the Broad/Orange processed mirrors, but report the Schliep version because it
gave the closest match to the GaGa empirical figures.  For MAQC, we downloaded
the public NCBI GEO GSE5350 series-matrix and annotation files:
\texttt{GSE5350-GPL570\_series\_matrix.txt.gz} for the Affymetrix arrays,
\texttt{GSE5350-GPL4097\_TaqMan\_series\_matrix.txt.gz} for the TaqMan qPCR
assays, and the corresponding GPL570/GPL4097 platform annotation files.  The
empirical results should therefore be read as a reproduction from currently
available public data, rather than a bit-for-bit rerun of the original GaGa
data objects.}  First, for the Armstrong
leukemia data comparing ALL and MLL samples, we used the recovered public
Schliep-filtered data set with 2194 genes and 42 samples.  The original GaGa
reproduction selected 360 genes at posterior FDR 0.05.  Replacing only the
posterior scorer by P-IG MCMC again selected 360 genes; 354 genes were common
to both lists, giving Jaccard overlap 0.967.  The Spearman correlation between
the original GaGa differential-expression score and the P-IG MCMC score was
0.999, with mean absolute score difference 0.0062.

Second, for the MAQC validation study, we fitted the GaGa empirical-Bayes
hyperparameters on all 54,675 first-site Affymetrix probes and then evaluated
the P-IG scorer on the 1893 Affymetrix probes mapped to qPCR assays.  This
mirrors the original empirical-Bayes use of all genes while keeping the P-IG
computation focused on the validation set.  The original GaGa validation AUC
was 0.794, MiGaGa2 gave 0.791, and limma gave 0.807 in our reproduction.  On
the same qPCR-mapped probes, the P-IG MCMC scorer gave AUC 0.739.  The P-IG
score remained highly correlated with the original GaGa score
(Spearman correlation 0.976), but was more conservative: at posterior FDR 0.05
the original GaGa scorer selected 1531 probes, whereas P-IG MCMC selected 1469.
The most probable pattern counts also shifted toward the null and intermediate
patterns: among the 1893 qPCR-mapped probes, original GaGa assigned 445, 71,
263, 106, and 1008 probes to patterns 0--4, respectively, while P-IG MCMC
assigned 498, 66, 414, 111, and 804 probes.

\begin{table}[t]
\centering
\caption{Empirical comparison between original GaGa scoring and P-IG MCMC
scoring.  Hyperparameters are estimated by the original GaGa empirical-Bayes
EM algorithm in both methods; P-IG MCMC changes only the posterior scoring
step.}
\label{tab:gaga-empirical}
\smallskip
\renewcommand{\arraystretch}{1.12}
\setlength{\tabcolsep}{5pt}
{\footnotesize
\begin{tabular}{@{}llrrrr@{}}
\toprule
\textbf{Data set} & \textbf{Method} & \textbf{Selected} &
\textbf{AUC} & \textbf{Spearman} & \textbf{Notes} \\
\midrule
Armstrong ALL--MLL & Original GaGa & 360 & -- & 1.000 & 2194 genes \\
                   & P-IG MCMC & 360 & -- & 0.999 & 354 overlapping calls \\
\midrule
MAQC validation & Original GaGa & 1531 & 0.794 & 1.000 & 1893 qPCR-mapped probes \\
                & P-IG MCMC & 1469 & 0.739 & 0.976 & Same EM hyperparameters \\
                & limma & -- & 0.807 & 0.987 & Reference method \\
\bottomrule
\end{tabular}}
\end{table}

These experiments clarify the role of the P-IG replacement.  The method is not
intended to change the empirical-Bayes GaGa hierarchy or to introduce a new
hyperparameter estimator.  Instead, it replaces the Stirling-based gamma-shape
posterior scoring step by an exact data-augmented MCMC calculation.  When the
shape parameters lie in the high-shape regimes typical of the Armstrong data,
the original approximation is already accurate and P-IG MCMC produces nearly
identical calls.  In low-shape, high-CV regimes, however, the oracle simulation
shows that the Stirling scorer can substantially distort posterior null
probabilities and posterior FDR decisions, while P-IG MCMC remains close to
exact quadrature.  The MAQC results suggest that on real validation data the
P-IG scorer can be somewhat more conservative than the original GaGa scorer.
Thus the main empirical message is one of robustness: P-IG MCMC preserves the
behavior of GaGa in ordinary regimes and provides a principled safeguard in
the regimes where the Stirling approximation is least reliable.
